\section{Giới thiệu nhóm và phạm vi đồ án}
\subsection{Bối cảnh và mục tiêu}
FloodRoute HCMC là đồ án Lab 1 của nhóm em trong môn Nhập môn Trí tuệ Nhân tạo. Mục tiêu của nhóm là mô hình hóa bài toán tìm đường giao hàng trên graph có hướng, cài đặt các thuật toán tìm kiếm và so sánh kết quả theo composite cost. Nhóm em dùng UCS làm mốc để kiểm tra A* và dùng các scenario traffic/flood để minh họa ảnh hưởng của môi trường.

\subsection{Thành viên và phân công}
\begin{tabularx}{\textwidth}{l l X}
\toprule
Thành viên & Mã số & Phạm vi chính\\
\midrule
Ngô Quang Thắng & 23127473 & Mô hình graph, cost engine, tích hợp backend\\
Vũ Lê Trọng Văn & 20127095 & Các thuật toán tìm kiếm và trace frontier\\
Nguyễn Duy Khang & 23127202 & Dataset, scenario và kiểm thử\\
Lưu Huy Minh Quang & 23127016 & Frontend, bản đồ và báo cáo\\
\bottomrule
\end{tabularx}

\subsection{Phạm vi}
Trong báo cáo, nhóm em sử dụng bản đồ chính gồm 65 landmark và 178 directed road-path edges. Các scenario peak/rain trong phần kiểm chứng dùng graph nhỏ có nhãn \texttt{SIMULATED}, không đại diện cho traffic thời gian thực.
